% Section 4.1, from lectures/L04/appendix/a_feedback.md.

\section{Feedback, and what one number decides}
\label{c4:app:a}

Chapter~\ref{c3:ch:filters} asserted two rules and used them. This section derives them, and the derivation hands back an error term that says how far from true they are.

\subsection{The loop, and the one number}
\label{c4:sec:a1}
An amplifier of gain $A$, a fraction $\beta$ of its output fed back in opposition:

\[
  A_{CL} = \frac{A}{1 + A\beta}
\]

$A\beta$ is the \textbf{loop gain}, written $T$: what a signal is multiplied by on one trip around the loop, and the only quantity in this section that matters. When $T$ is large,

\[
  A_{CL} \approx \frac{1}{\beta}
\]

\textbf{which depends on the feedback network alone, and that is the entire reason feedback is used.} $\beta$ is two resistors; $A$ varies by three between devices and drifts with temperature.

\textbf{The two rules of Chapter~\ref{c3:ch:filters} are this restated.} Gain set by $\beta$ alone is rule 2, the virtual short. Rule 1 is the amplifier's own property, not a feedback result at all.

\subsection{The error, which is the useful form}
\label{c4:sec:a2}
\[
  A_{CL} = \frac{1}{\beta} \cdot \frac{T}{1 + T} = \frac{1}{\beta}\left(1 - \frac{1}{1+T}\right)
\]

\textbf{The gain falls short of ideal by one part in $1 + T$,} which turns a judgement into arithmetic.

\bookfigure{lectures/L04/appendix/images/gain_error.png}{Gain error against open-loop gain on logarithmic axes, for closed-loop gains of ten, one hundred and one thousand. Each curve falls as one over the loop gain, and reaching 0.01 per cent error needs an open-loop gain of ten thousand times the closed-loop gain.}{c4:fig:gainerror}

\begin{narrowtable}{@{}ll@{}}
  \thead{Loop gain $T$} & \thead{Gain error} \\
  \midrule
  10 & 9 per cent \\
  100 & 1 per cent \\
  $10^3$ & 0.1 per cent \\
  $10^4$ & 0.01 per cent \\
\end{narrowtable}

An amplifier with $A = 10^5$ at a closed-loop gain of 10 has $\beta = 0.1$, so $T = 10^4$ and the error is 0.01 per cent.

At a closed-loop gain of 1000 it has $T = 100$ and 1 per cent error, usually unacceptable. \textbf{High closed-loop gain is expensive,} which the closed form does not show.

\subsection{What else $1 + T$ divides}
\label{c4:sec:a3}
\begin{itemize}
  \item \textbf{Distortion is divided by $1 + T$.} The loop sees the amplifier's own nonlinearity as an error to correct: 1 per cent open loop becomes 0.01 per cent closed.
  \item \textbf{Output impedance is divided by $1 + T$} for voltage feedback: 100 ohm open loop with $T = 10^4$ presents 10 milliohm.
  \item \textbf{Input impedance is multiplied by $1 + T$} for series feedback, which the non-inverting configuration uses. Inverting uses shunt feedback, which divides it: hence its bare $R_{in}$.
  \item \textbf{Gain sensitivity is divided by $1 + T$.} A 50 per cent change in $A$ moves $A_{CL}$ by 50 per cent divided by $1 + T$.
\end{itemize}
\textbf{What feedback does not fix.} More feedback cannot drive the remainder to zero, because $T$ itself falls with frequency, and nothing here touches input noise, offset or clipping. \textbf{A clipped amplifier has no loop gain at all, so feedback stops working exactly when it would be most useful.}

\subsection{What it costs: gain-bandwidth}
\label{c4:sec:a4}
Open-loop gain falls as a single pole from a few hertz upwards, at 20 dB per decade, so the product of gain and frequency is a constant:

\[
  \text{GBW} = A_{CL} \times f_{-3\ \text{dB}}
\]

A 1 MHz gain-bandwidth product gives 100 kHz at a closed-loop gain of 10, and 10 kHz at 100.

\textbf{Gain and bandwidth trade one for one,} which is why a high-gain stage is usually two stages of moderate gain: two stages of 10 have about six and a half times the bandwidth of one stage of 100, ten per stage less the 0.644 a cascade of coincident poles costs (\secref{c4:sec:a6}).

$T$ falls with the open-loop gain, so every benefit in \secref{c4:sec:a3} evaporates as frequency rises: 0.01 per cent gain error at DC is 1 per cent at a hundredth of the bandwidth. \textbf{Closed-loop distortion rises with frequency even when the amplifier's own does not.}

\subsection{Where feedback stops working}
\label{c4:sec:a5}
Each pole around the loop contributes up to 90 degrees of lag. \textbf{At 180 degrees negative feedback adds instead of subtracting, and if $T$ is still above one there, the circuit oscillates.}

One pole can never reach 180 degrees, so a single-pole amplifier with resistive feedback is unconditionally stable; two reach it only asymptotically. \textbf{Three reach it with gain to spare,} which is why a three-stage amplifier needs deliberate compensation: Chapter~\ref{c10:ch:amplifier}'s Miller capacitor.

Phase margin, compensation strategy and the Nyquist criterion are a course of their own.

\subsection{Active filters}
\label{c4:sec:a6}
An amplifier can do better than isolate two RC sections: it can put feedback around the passive network and produce a response no passive RC network can.

A \textbf{Sallen-Key} low-pass is two resistors, two capacitors and one amplifier as a follower, with the first capacitor returned to the output rather than to ground. That capacitor sees the output, feeds energy back into the network, and \textbf{the pole pair becomes complex.}

A passive RC cascade makes only real poles, and real poles give a soft corner: two coincident are 6 dB down. A complex pair can be 3 dB down, or peaked, and far sharper.

\begin{booktable}{@{}LlL@{}}
  \thead{Realisation} & \thead{Poles} & \thead{3 dB point of two sections} \\
  \midrule
  Two RC sections cascaded directly & Real, split apart & 0.374 of one section's corner \\
  Two RC sections with a buffer between & Real, coincident & 0.644 \\
  Sallen-Key, Q = 0.707 & Complex pair & 1.0, and a much sharper transition \\
\end{booktable}

\textbf{One part buys a filter the passive components cannot make at any value.} That is a better argument for an op-amp than gain is, and the argument Chapter~\ref{c8:ch:follower} makes for an emitter follower.

\subsection{What this section is blind to}
\label{c4:sec:a7}
\begin{itemize}
  \item \textbf{Stability, properly.} The condition is stated in \secref{c4:sec:a5} and never developed. A design that needs a phase margin computed needs a different course.
  \item \textbf{Feedback topologies.} There are four, by whether voltage or current is sensed and fed back. This section assumes voltage sensing and describes two.
  \item \textbf{Noise.} Feedback does not reduce noise generated at the amplifier's input, which in most low-noise designs decides everything.
\end{itemize}
